\documentclass{article}

\usepackage{booktabs}
\usepackage{tabularx}
\usepackage{float}

\title{Development Plan\\\progname}

\author{\authname}

\date{}

\input{../Comments.text}
\input{../Common.text}

\begin{document}

\maketitle

\begin{table}[hp]
\caption{Revision History} \label{TblRevisionHistory}
\begin{tabularx}{\textwidth}{llX}
\toprule
\textbf{Date} & \textbf{Developer(s)} & \textbf{Change}\\
\midrule
Date1 & Name(s) & Description of changes\\
Date2 & Name(s) & Description of changes\\
... & ... & ...\\
\bottomrule
\end{tabularx}
\end{table}

\newpage{}

\wss{Put your introductory blurb here.  Often the blurb is a brief roadmap of
what is contained in the report.}

\wss{Additional information on the development plan can be found in the
\href{https://gitlab.cas.mcmaster.ca/courses/capstone/-/blob/main/Lectures/L02b_POCAndDevPlan/POCAndDevPlan.pdf?ref_type=heads}
{lecture slides}.}

\section{Confidential Information?}

\wss{State whether your project has confidential information from industry, or
not.  If there is confidential information, point to the agreement you have in
place.}

\wss{For most teams this section will just state that there is no confidential
information to protect.}
\section{IP to Protect}

\wss{State whether there is IP to protect.  If there is, point to the agreement.
All students who are working on a project that requires an IP agreement are also
required to sign the ``Intellectual Property Guide Acknowledgement.''}

\section{Copyright License}

\wss{What copyright license is your team adopting.  Point to the license in your
repo.}

\section{Team Meeting Plan}

\wss{This section is for meetings between the team members, and meetings between 
the team and their supervisor/advisor (as appropriate).}

\wss{How often will the team meet? where? when?}

\wss{If the meeting is a physical location (not virtual), out of an abundance of
caution for safety reasons you shouldn't put the location online}

\wss{To help schedule meetings, you can use the lecture/tutorial time when we do 
not have anything else scheduled.}

\wss{How often will you meet with your supervisor/advisor?  when?  where?}

\wss{Will meetings be virtual?  At least some meetings should likely be
in-person.}

\wss{How will the meetings be structured?  There should be a chair for all meetings.  
There should be an agenda for all meetings.}

\section{Team Communication Plan}

\wss{Issues on GitHub should be part of your communication plan.}

\section{Team Member Roles}

\wss{You should identify the types of roles you anticipate, like notetaker,
leader, meeting chair, reviewer.  Assigning specific people to those roles is
not necessary at this stage.  In a student team the role of the individuals will
likely change throughout the year.}

\section{Workflow Plan}

\begin{itemize}
	\item How will you be using git, including branches, pull request, etc.?
	\item How will you be managing issues, including template issues, issue
	classification, etc.?
  \item Use of CI/CD
\end{itemize}

This section describes how the team will use Git and GitHub to collaborate on both source code and documentation, how work will be tracked through issues, and how Continuous Integration and Continuous Deployment (CI/CD) will keep the project stable. 
It is written to serve as a guide for current and future contributors. The repository is available at \href{https://github.com/davislenover/BurnNEarn}{BurnNEarn GitHub Repository}.

\subsection{Branches}

The repository will have exactly one permanent branch:

\begin{itemize}
  \item \texttt{main}: the protected, stable branch representing the latest reviewed state of the project. Direct pushes are not permitted. 
  All changes must arrive through a pull request that has been approved by at least one other team member and has passed all CI checks.
\end{itemize}

All other branches are short-lived and must follow the naming convention \texttt{<type>/<name>-<description>}, where \texttt{type} is one of the categories in Table~\ref{TblBranchTypes}, \texttt{name} is the author's first name in lowercase, and \texttt{description} is a brief summary of the work in two to four words, separated by hyphens.

\begin{table}[hp]
\caption{Branch Types and Naming Examples} \label{TblBranchTypes}
\begin{tabularx}{\textwidth}{lXl}
\toprule
\textbf{Type} & \textbf{Used For} & \textbf{Example}\\
\midrule
\texttt{docs} & Documentation sections & \texttt{docs/farid-dev-plan-workflow}\\
\texttt{feature} & New app functionality & \texttt{feature/davis-step-sync}\\
\texttt{fix} & Bug fixes & \texttt{fix/angad-portfolio-rounding}\\
\texttt{chore} & CI, configuration, and repository maintenance & \texttt{chore/preston-latex-build}\\
\bottomrule
\end{tabularx}
\end{table}

Branches are not shared between members; each member creates their own branch for each unit of work. 
For documentation, each member creates one branch per deliverable covering the sections assigned to them. For code, each feature or fix gets its own branch. Once a branch's pull request is merged into \texttt{main}, the branch is deleted. Before opening a pull request, the author updates their branch with the latest changes from \texttt{main} and resolves any merge conflicts locally, so that conflicts are never resolved during review.

\subsection{Commits}

Commits should be small and focused on a single change, rather than combining unrelated edits. 
Commit messages should begin with a short summary in the imperative mood (e.g., ``Add step count permission request'') and reference the related issue number where applicable (e.g., \texttt{\#14}), which links the commit to that issue on GitHub. 
Issues are closed through the pull request rather than individual commits, as described in Section~\ref{SecPullRequests}.

\subsection{Pull Requests} \label{SecPullRequests}

When the work on a branch is complete, the author follows these steps:

\begin{enumerate}
  \item Open a pull request into \texttt{main} with a title summarizing the change and a description of what was changed and why. The description must link the related issue using a closing keyword (e.g., \texttt{Closes \#14}), which automatically closes the issue when the pull request is merged.
  \item Request at least one reviewer. The reviewer should be a member working in a related area, or the whole team if the change affects everyone, such as shared document sections or the overall architecture.
  \item Address the reviewer's comments, which are left directly on the relevant lines of code or documentation.
  \item Once the pull request has one approval and all CI checks pass, the author merges it using a \textit{merge commit}, which preserves every individual commit from the branch on \texttt{main}. This keeps a complete record of each member's contributions throughout the project. The author then deletes the branch.
\end{enumerate}

Reviewers are expected to respond to a review request within 48 hours, or within 24 hours during the final three days before a deadline. 
For documentation, reviewers also check the changes against the deliverable's checklist and rubric before approving.

\subsection{Tags}

Deliverables are marked directly from the \texttt{main} branch, so a Git tag will be created on \texttt{main} at each deliverable's deadline to record exactly what was marked. 
Document deadlines use the format \texttt{<document>-<revision>} (e.g., \texttt{dev-plan-rev0}, \texttt{srs-rev0}, \texttt{design-rev-1}), and demonstrations use the format \texttt{<demo>-demo} (e.g., \texttt{poc-demo}, \texttt{rev0-demo}, \texttt{final-demo}). 
The team member who merges the final pull request before a deadline is responsible for creating its tag. 
Tags allow the team to return to the exact state of the project at each deadline, and to compare revisions when addressing feedback.

\subsection{Issue Management}

GitHub Issues will serve as the team's ticketing system, so that all work is visible, assigned, and traceable.

\begin{itemize}
  \item \textbf{Issue size:} Each issue should represent work that can be completed in a single sitting of roughly two to three hours. Larger tasks are split into multiple issues. For documentation, every section of each document has its own issue, titled with its deliverable and section (e.g., \texttt{[Development Plan] 7 Workflow Plan}).
  \item \textbf{Ownership:} Every task issue is assigned to one owner, who is responsible for completing it and moving it through the project board described in Section~\ref{SecGitHubProjects}.
  \item \textbf{Milestones:} Each deliverable is a milestone whose due date matches the deadline in Table~\ref{TblSchedule}. Every issue related to a deliverable is added to its milestone, so the team can see how much work remains before each deadline.
  \item \textbf{Meetings and lectures:} An issue is created for every team meeting, supervisor meeting, and lecture, and is tracked in its own column on the project board. At least 24 hours before a meeting, the agenda and any questions are added to the issue. Afterwards, notes and action items are recorded, each action item becomes its own issue, and the meeting issue is closed. Lecture issues record the topic, date and time, and attendance.
\end{itemize}

The repository includes the issue templates listed in Table~\ref{TblIssueTemplates}, so that every issue contains consistent information.

\begin{table}[H]
  \caption{Issue Templates} \label{TblIssueTemplates}
  \begin{tabularx}{\textwidth}{lX}
    \toprule
    \textbf{Template} & \textbf{Contents}\\
    \midrule
      Documentation Task & Deliverable, section, and the checklist items the section must satisfy\\
      Feature & Description of the functionality and acceptance criteria\\
      Bug Report & Steps to reproduce, expected and actual behaviour, and device and operating system\\
      Meeting & Meeting type, date and time, attendance, agenda, questions, notes, and action items\\
      Lecture & Topic, date and time, and attendance\\
    \bottomrule
  \end{tabularx}
\end{table}

Issues are classified using the labels in Table~\ref{TblLabels}. Every issue that belongs to a deliverable receives that deliverable's label, and every issue receives one type label and, where relevant, one area label. 
Together, these allow the team to filter issues by deliverable, by kind of work, or by part of the app.

\begin{table}[H]
  \caption{Issue Labels} \label{TblLabels}
  \begin{tabularx}{\textwidth}{llX}
    \toprule
    \textbf{Category} & \textbf{Label} & \textbf{Meaning}\\
    \midrule
      % Deliverable & \texttt{Problem Statement \& Goals} & Problem Statement and Goals document\\
      Deliverable & \begin{tabular}[t]{@{}l@{}}\texttt{Problem Statement}\\\texttt{and Goals}\end{tabular} & Problem Statement and Goals document\\
      & \texttt{Development Plan} & Development Plan and Proof of Concept Plan\\
      & \texttt{Requirements} & Software Requirements Specification\\
      & \texttt{Hazard Analysis} & Hazard Analysis document\\
      & \texttt{V\&V Plan} & Verification and Validation Plan\\
      & \texttt{Design Document} & Design Document (all revisions)\\
      & \texttt{V\&V Report} & Verification and Validation Report\\
      & \texttt{Demo} & Proof of Concept, Revision 0, Final, and EXPO demonstrations\\
      \midrule
      Type & \texttt{documentation} & Writing or revising a document section\\
      & \texttt{reflection} & Reflection components of a deliverable\\
      & \texttt{feature} & New functionality\\
      & \texttt{enhancement} & Improvement to existing functionality\\
      & \texttt{bug} & A defect that must be fixed\\
      & \texttt{meeting} & Team or supervisor meeting\\
      & \texttt{lecture} & Course lecture\\
      \midrule
      Area & \texttt{frontend} & User interface of the mobile app\\
      & \texttt{backend} & Server, database, and accounts\\
      & \texttt{health-data} & Fitness tracking and health platform integration\\
      & \texttt{market-engine} & Asset prices, portfolios, and market scenarios\\
      & \texttt{social} & Social feed and sharing\\
    \bottomrule
  \end{tabularx}
\end{table}

\subsection{Continuous Integration and Continuous Deployment}


\section{Project Decomposition and Scheduling}
\subsection{GitHub Projects} \label{SecGitHubProjects}

The team will use a GitHub Project board to track all issues. The board will have the columns \textit{Backlog}, \textit{In Progress}, \textit{In Review}, and \textit{Done}. Issues move to \textit{In Review} when a pull request is opened, and to \textit{Done} automatically when the linked issue is closed by a merged pull request. The board can be filtered by milestone, assignee, or label, giving the team and supervisor a clear view of who is working on what and what remains before each deadline.

The GitHub Project can be found here: \href{https://github.com/users/davislenover/projects/4/views/1}{FitFolio GitHub Project}.

\subsection{Schedule}

The project will be scheduled around the course deliverables, with each deliverable set as a GitHub milestone. The major deadlines, taken from the course outline, are:

\begin{table}[H] \label{TblSchedule}
\centering
\renewcommand{\arraystretch}{1.3}
\begin{tabular}{l c}
\toprule
\textbf{Deliverable} & \textbf{Week} \\
\midrule
Team Formed, Project Selected & Week 03 \\
Problem Statement, POC Plan, Development Plan  & Week 04 \\
Req. Doc. and Hazard Analysis Revision 0 & Week 06 \\
V\&V Plan Revision 0 &  Week 08 \\
Design Document Rev -1 and CD-Fall Rev 0 & Week 10 \\
Proof of Concept Demo + Interviews &  Weeks 11--12 \\
Design Document Revision 0 &  Week 16 \\
Revision 0 Demo + Interviews &  Weeks 18--19 \\
V\&V Report Rev 0 and CD-Winter Rev 0 & Week 22 \\
Final Demonstration (Revision 1) & Week 24 \\
EXPO Demonstration & Week 26 \\
Final Documentation (Revision 1) & Week 26 \\
\bottomrule
\end{tabular}
\caption{Major project deadlines from the course outline} \label{TblSchedule}
\label{tab:schedule}
\end{table}
\newpage{}

\begin{itemize}
  \item How will you be using GitHub projects?
  \item Include a link to your GitHub project
\end{itemize}

\wss{How will the project be scheduled?  This is the big picture schedule, not
details. You will need to reproduce information that is in the course outline
for deadlines.}

\section{Proof of Concept Demonstration Plan}

What is the main risk, or risks, for the success of your project?  What will you
demonstrate during your proof of concept demonstration to convince yourself that
you will be able to overcome this risk?

\section{Expected Technology}

\wss{What programming language or languages do you expect to use?  What external
libraries?  What frameworks?  What technologies.  Are there major components of
the implementation that you expect you will implement, despite the existence of
libraries that provide the required functionality.  For projects with machine
learning, will you use pre-trained models, or be training your own model?  }

\wss{The implementation decisions can, and likely will, change over the course
of the project.  The initial documentation should be written in an abstract way;
it should be agnostic of the implementation choices, unless the implementation
choices are project constraints.  However, recording our initial thoughts on
implementation helps understand the challenge level and feasibility of a
project.  It may also help with early identification of areas where project
members will need to augment their training.}

Topics to discuss include the following:

\begin{itemize}
\item Specific programming language
\item Specific libraries
\item Pre-trained models
\item Specific linter tool (if appropriate)
\item Specific unit testing framework
\item Investigation of code coverage measuring tools
\item Specific plans for Continuous Integration (CI), or an explanation that CI
  is not being done
\item Specific performance measuring tools (like Valgrind), if
  appropriate
\item Tools you will likely be using?
\end{itemize}

\wss{git, GitHub and GitHub projects should be part of your technology.}

\section{Use of AI and LLMs}

\wss{This section is a continuation of the previous section.  The role of LLMs
in your project as important enough to warrant their own section.  
How will your team be using LLMs for code and documentation development?  
What tools will you be using?  How will you verify the output of your LLMs?}

\section{Coding Standard}

\wss{What coding standard will you adopt?}

\newpage{}

\section{Appendix --- Generative AI Use Disclosure}

\wss{ChatGPT (version 5) was used to brainstorm ideas for the template for this disclose section.}

\wss{This section is about the use of AI to write this document, not your
planned use of AI for your project.  You discuss that topic in one of the
sections in the body of this report.}

\subsection{AI Tools Used}

\wss{Give the name of the AI tool(s) used and the version.}

\subsection{How and Where Used}

\wss{Explain what the tool(s) were used for.  Examples include: brainstorming,
explaining a technical concept, reviewing the document, generating or modifying
text, generating or modifying diagrams, identifying errors or omissions or
inconsistencies.}

\wss{For each of the uses, identify which parts of the document were affected.}

\wss{You do not need to report every interaction, but you should disclose the 
uses that materially contributed to your work}

\subsection{Verification of AI-Generated Content}

\wss{\begin{itemize}
    \item Describe how the team checked AI-generated or AI-assisted material
    before including it in the document.
    \item In particular, verify factual claims, technical assertions,
    requirements, calculations, references, and other information for which an
    AI system may produce incorrect or fabricated results.
    \item \textbf{The team is responsible for the accuracy, completeness,
consistency, and appropriateness of the submitted document, regardless of
whether material was generated by a human or an AI system.}
\end{itemize}}

\newpage{}

\section*{Appendix --- Reflection}

\wss{Not required for CAS 741}

\input{../Reflection.text}

\begin{enumerate}
    \item Why is it important to create a development plan prior to starting the
    project?
    \item In your opinion, what are the advantages and disadvantages of using
    CI/CD?
    \item What disagreements did your group have in this deliverable, if any,
    and how did you resolve them?
\end{enumerate}

\newpage{}

\section*{Appendix --- Team Charter}

\wss{borrows from
\href{https://engineering.up.edu/industry_partnerships/files/team-charter.pdf}
{University of Portland Team Charter}}

\subsection*{External Goals}

\wss{What are your team's external goals for this project? These are not the
goals related to the functionality or quality fo the project.  These are the
goals on what the team wishes to achieve with the project.  Potential goals are
to win a prize at the Capstone EXPO, or to have something to talk about in
interviews, or to get an A+, etc.}

\subsection*{Attendance}

\subsubsection*{Expectations}

\wss{What are your team's expectations regarding meeting attendance (being on
time, leaving early, missing meetings, etc.)?}

\subsubsection*{Acceptable Excuse}

\wss{What constitutes an acceptable excuse for missing a meeting or a deadline?
What types of excuses will not be considered acceptable?}

\subsubsection*{In Case of Emergency}

\wss{What process will team members follow if they have an emergency and cannot
attend a team meeting or complete their individual work promised for a team
deliverable?}

\subsection*{Accountability and Teamwork}

\subsubsection*{Quality} 

\wss{What are your team's expectations regarding the quality
of team members' preparation for team meetings and the quality of the
deliverables that members bring to the team?}

\subsubsection*{Attitude}

\wss{What are your team's expectations regarding team members' ideas,
interactions with the team, cooperation, attitudes, and anything else regarding
team member contributions?  Do you want to introduce a code of conduct?  Do you
want a conflict resolution plan?  Can adopt existing codes of conduct.}

\subsubsection*{Stay on Track}

\wss{What methods will be used to keep the team on track? How will your team
ensure that members contribute as expected to the team and that the team
performs as expected? How will your team reward members who do well and manage
members whose performance is below expectations?  What are the consequences for
someone not contributing their fair share?}

\wss{You may wish to use the project management metrics collected for the TA and
instructor for this.}

\wss{You can set target metrics for attendance, commits, etc.  What are the
consequences if someone doesn't hit their targets?  Do they need to bring the
coffee to the next team meeting?  Does the team need to make an appointment with
their TA, or the instructor?  Are there incentives for reaching targets early?}

\subsubsection*{Team Building}

\wss{How will you build team cohesion (fun time, group rituals, etc.)? }

\subsubsection*{Decision Making} 

\wss{How will you make decisions in your group? Consensus?  Vote? How will you
handle disagreements? }

\end{document}